\clearpage
\VerdierPhysicalPageAnchor{0175}%
\begingroup\small\setlength{\parskip}{0pt}\setlength{\emergencystretch}{1.5em}
\setlength{\abovedisplayskip}{4pt plus 2pt minus 2pt}%
\setlength{\belowdisplayskip}{4pt plus 2pt minus 2pt}%

\VerdierSemanticAnchor{ch3d-1-1-2-cont}
$K(\mathcal A)$ における完全三角の定義（第 I 章，3.3.1）と標準コホモロジー的
関手 $H\colon K(\mathcal A)\longrightarrow\mathcal A$ の性質から、圏
$K^{*,*'}(\mathcal A)\subset K(\mathcal A)$ は $K(\mathcal A)$ の充満三角
部分圏（第 II 章，2.1.5；$*={}$「空」のときは厳密充満）であり、飽和部分圏
（第 II 章，2.1.6）であることが直ちに従う。

\VerdierSemanticAnchor{ch3d-1-1-3}
\medskip\noindent\textbf{1.1.3.}\quad
$\mathcal A$ の複体 $X$ は、すべての整数 $i$ に対して $H^i(X)=0$ であるとき
\emph{非輪状}であるという。$\operatorname{Ac}^{*}(\mathcal A)$ を、
$K^{*}(\mathcal A)$ の対象である非輪状複体で定まる $K^{*}(\mathcal A)$ の
充満部分圏とする。$K(\mathcal A)$（それぞれ
$\operatorname{comp}(\mathcal A)$）の射 $X\xrightarrow{u}Y$ は、すべての
整数 $i$ に対して射 $H^i(X)\xrightarrow{H^i(u)}H^i(Y)$ が同型であるとき
\emph{擬同型}と呼ばれる。$\operatorname{Qis}^{*}(\mathcal A)$
（混同のおそれがないときは $\operatorname{Qis}^{*}$）を、
$K^{*}(\mathcal A)$ の擬同型全体の集合とする。
$K^{*}(\mathcal A)$ の部分圏 $\operatorname{Ac}^{*}(\mathcal A)$ は、
$K^{*}(\mathcal A)$ の厳密充満かつ飽和な三角部分圏である。
$\operatorname{Qis}^{*}(\mathcal A)$
は $K^{*}(\mathcal A)$ の射からなる飽和乗法的系である（第 II 章，2.1.4）。
$K^{*}(\mathcal A)$ の完全三角 $(X,Y,Z,u,v,w)$ において、射 $u$ が擬同型で
あるための必要十分条件は複体 $Z$ が非輪状であることである。言い換えれば
（第 II 章，2.1.15）、$\operatorname{Qis}^{*}(\mathcal A)$ は、充満三角
部分圏 $\operatorname{Ac}^{*}(\mathcal A)$ に付随する $K^{*}(\mathcal A)$ の
乗法的系である。

\VerdierSemanticAnchor{ch3d-1-1-4}
\medskip\noindent\textbf{1.1.4.}\quad
$M$ を $\mathcal A$ の対象の集合とする。$\mathcal A$ の複体は、そのすべての
成分が $M$ の元であるとき\emph{$M$ 型}であるという。
$K^{*}(\mathcal A)$ または $\operatorname{comp}^{*}(\mathcal A)$ の射
$u\colon X\longrightarrow Y$ は、$u$ が擬同型であり、かつ $Y$（それぞれ
$X$）が $M$ 型複体であるとき、$X$（それぞれ $Y$）の\emph{$M$ 型右分解}
（それぞれ\emph{左分解}）と呼ばれる。多くの場合、文脈上混同のおそれが
なければ、「右」（それぞれ「左」）を省く。また多くの場合、同じく言葉を
濫用して、対象 $Y$ の右分解（それぞれ左分解）の終域（それぞれ始域）を
$Y$ の右分解（それぞれ左分解）と呼ぶ。
$\operatorname{comp}_{\mathcal A}^{*}(M)$、または単に
$\operatorname{comp}^{*}(M)$（それぞれ $K_{\mathcal A}^{*}(M)$、または単に
$K^{*}(M)$）を、$M$ 型複体で定まる
$\operatorname{comp}^{*}(\mathcal A)$（それぞれ $K^{*}(\mathcal A)$）の
充満部分圏とする。

\VerdierSemanticAnchor{ch3d-1-1-5}
\medskip\noindent\textbf{1.1.5.}\quad
$M$ が圏 $\mathcal A$ の射影対象（それぞれ入射対象）全体の集合であるとき、
その成分が射影対象（それぞれ入射対象）である複体を、上の規約に従って
\emph{射影型複体}（それぞれ\emph{入射型複体}）と呼ぶ。射影型複体
（それぞれ入射型複体）と、圏 $\operatorname{comp}(\mathcal A)$ の射影対象
（それぞれ入射対象）を混同してはならない。
$\operatorname{comp}(\mathcal A)$ の対象が射影対象（それぞれ入射対象）で
あるための必要十分条件は、それが射影型（それぞれ入射型）であり、かつ
零複体にホモトピックであることである。
\endgroup
